
The data reported here were collected from December 23, 2021 to May 11, 2022, under stable detector conditions. The cathode and gate electrodes \cite{LINEHANgrids} established a drift field of \SI{193}{V/cm}, determined by electrostatic simulation to vary by \SI{4}{\percent} over the volume considered in this analysis. The gate and anode electrodes established a gas extraction field of \SI{7.3}{kV/cm} at radial position $r=0$. Twelve TPC and two skin PMTs, with no specific position correlation, developed malfunctioning connections or excessive noise during commissioning and were disabled prior to the run. The temperature and pressure of the LXe were stable to within 0.2$\%$, at \SI{174.1}{\kelvin} (at the TPC bottom) and \SI{1.791}{bar(a)}. The liquid level was stable to within \SI{10}{\micro\meter}, measured by precision capacitance sensors. The full xenon complement of \SI{10}{t} was continuously purified at \SI{3.3}{t/day} through a hot getter system, and the observed electron lifetime against attachment on electronegative impurities was 
between \SI{5000}{\micro\s} and \SI{8000}{\micro\s}, much longer than the \SI{951}{\micro\s} maximum drift time in the TPC. 

The data acquisition (DAQ) system records events triggered by a digital filter sensitive to S2 signals in the TPC, reaching full efficiency for S2 pulses with six extracted electrons at a typical rate of \SI{5}{Hz}. A time window of \SI{2}{ms} before and \SI{2.5}{ms} after each trigger is recorded, constituting an event. Zero-suppressed waveforms from all PMT channels, including low- and high-gain amplification paths for TPC and OD PMTs, are recorded for every trigger with single photoelectron efficiencies averaging \SI{94}{\percent}, \SI{86}{\percent}, and $>$\SI{95}{\percent} for the TPC, skin, and OD PMTs, respectively.

Event properties are reconstructed through analysis of the PMT waveform shapes, timings, and distributions. Raw waveform amplitudes are normalized by the PMT and amplifier gains and summed separately within the TPC, skin, and OD. Integrated waveform area is reported in photons detected (phd) at each PMT, accounting for the double photoelectron effect in response to vacuum ultraviolet photons~\cite{faham2015measurements, LOPEZPAREDESdpe}. Pulse boundaries are identified on the summed waveforms using filters tuned for prototypical pulse shapes in each detector. Pulses in the TPC are further classified as S1 or S2 based on their hit pattern and pulse shape. S1 pulses are required to have signals above the electronic noise threshold in at least three PMTs. The time ordering of the most prominent S1 and S2 pulses in each event is then used to identify single-scatter (one S1 preceding one S2) and multiscatter (one S1 preceding multiple S2s) events. The transverse $(x,y)$ location of events is determined by the PMT hit pattern of S2 light from the extracted electrons, using the \textsc{Mercury} algorithm~\cite{Mercury}. The algorithm was tuned using uniformly distributed radioactive sources in the TPC and has a 1$\sigma$ resolution of \SI{4}{mm} for S2 signals of \SI{3000}{phd}. The resolution worsens by approximately a factor of 2~near the TPC wall due to asymmetric light collection at the TPC edge. The location along the cylinder ($z$) axis is inferred from the drift time, and has a 1$\sigma$ resolution of \SI{0.7}{mm} for events near the cathode electrode.

LZ uses radioactive sources to correct for spatial variation in response across the TPC and to calibrate the detector response to ER and NR events. ER calibration events are obtained using dispersed sources \KrEightThreem\ and \XeOneThreeOnem\ before and during the WIMP search and tritiated methane (CH$_3$T) postsearch. The tritium source is important for understanding the response to low-energy ER events, the most prominent background component in the run. Localized NR calibration events are created using a deuterium-deuterium (DD) generator that emits monoenergetic \SI{2.45}{MeV} neutrons~\cite{adelphi, Verbus_others_2017, LUX_DD} along a conduit through the water tank at approximately 10 cm below the liquid surface and AmLi sources~\cite{MOZHAYEVAmLi} deployed between the walls of the cryostat vessels in three azimuthal positions and three $z$ positions, a total of nine positions.

Using the dispersed sources, the S1 signal is normalized to the geometric center of the detector, using a correction in $x$, $y$, and drift time; this normalized value is called \sonec. The S2 signal is normalized to a signal at the radial center and top (shortest drift time) of the detector; this normalized value is called \stwoc. The size of the S1 corrections is on average \SI{9}{\percent} and comes primarily from variations in light collection efficiency and PMT quantum efficiency. The size of the S2 corrections is on average \SI{11}{\percent} in the $(x, y)$ plane and comes primarily from nonoperational PMTs and extraction-field nonuniformity caused by electrostatic deflection of the gate and anode electrodes. The S2 correction in $z$ is due to electron attachment on impurities and averages \SI{7}{\percent}.  Corrected parameters are uniform across the TPC to within $3$\%. 

To reproduce the TPC response to ER and NR events, detector and xenon response parameters of the \textsc{nest} 2.3.7~\cite{szydagis_m_2022_6534007} ER model are tuned to match the median and widths of the \tritium\ calibration data in log$_{10}$\stwoc-\sonec\, space, and to match the reconstructed energies of the \KrEightThreem~(\SI{41.5}{keV}), \XeOneTwoNinem~(\SI{236}{keV}), and \XeOneThreeOnem~(\SI{164}{keV}) peaks. The photon detection efficiency $g_1$ is determined to be \SI{0.114(2)}{phd/photon} and the gain of the ionization channel $g_2$ to be \SI{47.1(11)}{phd/electron}~\footnote{See Supplemental Material at \url{\supplementalurl} for the full description of the \textsc{nest} parameters and a header file to configure the software for this model}. The \tritium\ data are best modeled with the \textsc{nest} recombination skewness model~\cite{akerib2020discrimination} disabled, and comparisons between the tuned model and tritium data using several statistical tests show consistency throughout the full tritium ER distribution~\cite{kolmogorov1933sulla,smirnov1939estimation, andersondarling, shapiro1965analysis}. The \textsc{nest}  ER model also includes effects from electron capture decays~\cite{temples2021measurement} when making predictions from electron capture background sources. The parameters of the ER model were propagated to the \textsc{nest} NR model and found to be in good agreement with  DD calibration data, matching NR band means and widths to better than \SI{1}{\percent} and \SI{4}{\percent} in log$_{10}$\stwoc, respectively. Further checks comparing DD and AmLi neutron calibrations agree to 1\%. Figure \ref{fig:S1S2} shows the \tritium\ and DD neutron data compared to the calibrated model. 

\begin{figure}[htbp]
  \centering
  \includegraphics[ width=\linewidth]{figures/SR1WS_calOnly_0629.pdf}
  \caption{Calibration events in $\log_{10}$\stwoc-\sonec\ for the tritium\ source (dark blue points, 5343 events) and the DD neutron source (orange points, 6324 events). Solid blue (red) lines indicate the median of the ER (NR) simulated distributions, and the dotted lines indicate the {\SI{10}{\percent}} and {\SI{90}{\percent}} quantiles. Thin gray lines show contours of constant electron-equivalent energy (keV$_{\text{ee}}$) and nuclear recoil energy (keV$_{\text{nr}}$).}
  \label{fig:S1S2}
\end{figure}
